\subsubsection{Effective planning horizon (post hoc)}\label{app:horizon}
% src: raea4/agentenv/h_horizon.py -> results/h_horizon_{core,x,m}.json|txt; scripts/make_u_outputs.py
At every logged decision of the K2e runs (calls and the decision to stop; the verdict is not compared) the belief is recomputed from the logged reports and the action compared with the $k$-step look-ahead policy $\pi_k$ (\cref{sec:setup}; $k=1,\dots,4,\infty$), either on the true model or, for the as-if planner of \cref{app:asif}, on a model with the replica's fault odds multiplied by 4 that both tracks the belief and plans; an action within $10^{-6}$ of the best depth-$k$ value agrees (\cref{tab:horizon}).
First decisions agree most with full look-ahead on inflated fault odds (first domain: terra \uHCoreTerraGFourFirstKInf\%, Sonnet \uHCoreSonnetGFourFirstKInf\%, against \uHCoreTerraGOneFirstKInf\% and \uHCoreSonnetGOneFirstKInf\% with the true odds).
No horizon describes the later decisions: with the true odds, agreement is at most \uHCoreTerraGOneLaterMax\% (terra) and \uHCoreSonnetGOneLaterMax\% (Sonnet) for any $k$, and in the fault-free matched tasks at most \uHMZTerraLaterMax\% and \uHMZSonnetLaterMax\%. Sonnet's later decisions agree somewhat more with one-step look-ahead (fault-free: \uHMZSonnetLaterKOne\% against \uHMZSonnetLaterKInf\% for $k=\infty$). The later-step shortfall is therefore not a shorter look-ahead of this exact kind, and no effective horizon is identified.
\begin{table}[!htbp]
\centering
\caption{Share of decisions (\%) agreeing with the $k$-step look-ahead policy $\pi_k$ on the true model ($g{=}1$) or with the replica's fault odds multiplied by $g{=}4$ (post hoc; ties agree). First: the first decision of an episode; later: all subsequent decisions up to and including the stop.}
\label{tab:horizon}
\scriptsize\setlength{\tabcolsep}{3pt}
% generated by scripts/make_u_outputs.py from raea4/agentenv/results/u_units_*.json, r_robust.json, h_horizon_*.json; do not edit
\begin{tabular}{@{}lllcrrrrr@{}}
\toprule
Set & Run & $g$ & Decisions ($n$) & $k{=}1$ & 2 & 3 & 4 & $\infty$ \\
\midrule
first domain (CORE) & flash & 1 & first (100) & 40 & 44 & 44 & 40 & 40 \\
first domain (CORE) & flash & 1 & later (335) & 39 & 41 & 43 & 43 & 42 \\
first domain (CORE) & flash & 4 & first (100) & 36 & 67 & 73 & 72 & 73 \\
first domain (CORE) & flash & 4 & later (335) & 35 & 37 & 38 & 38 & 38 \\
\addlinespace
first domain (CORE) & terra & 1 & first (100) & 65 & 65 & 65 & 61 & 61 \\
first domain (CORE) & terra & 1 & later (247) & 71 & 70 & 69 & 68 & 67 \\
first domain (CORE) & terra & 4 & first (100) & 54 & 80 & 84 & 84 & 86 \\
first domain (CORE) & terra & 4 & later (247) & 68 & 75 & 76 & 76 & 76 \\
\addlinespace
first domain (CORE) & Sonnet & 1 & first (100) & 62 & 63 & 63 & 59 & 59 \\
first domain (CORE) & Sonnet & 1 & later (209) & 69 & 60 & 59 & 58 & 58 \\
first domain (CORE) & Sonnet & 4 & first (100) & 51 & 80 & 85 & 87 & 89 \\
first domain (CORE) & Sonnet & 4 & later (209) & 73 & 72 & 70 & 69 & 70 \\
\addlinespace
fault-rate sweep & flash & 1 & first (200) & 21 & 37 & 38 & 42 & 42 \\
fault-rate sweep & flash & 1 & later (765) & 38 & 42 & 42 & 42 & 42 \\
fault-rate sweep & flash & 4 & first (200) & 26 & 54 & 56 & 64 & 64 \\
fault-rate sweep & flash & 4 & later (765) & 35 & 41 & 41 & 41 & 40 \\
\addlinespace
fault-rate sweep & terra & 1 & first (200) & 40 & 60 & 62 & 66 & 66 \\
fault-rate sweep & terra & 1 & later (599) & 64 & 69 & 70 & 71 & 71 \\
fault-rate sweep & terra & 4 & first (200) & 36 & 66 & 79 & 86 & 86 \\
fault-rate sweep & terra & 4 & later (599) & 66 & 75 & 77 & 78 & 79 \\
\addlinespace
fault-rate sweep & Sonnet & 1 & first (200) & 38 & 60 & 61 & 64 & 65 \\
fault-rate sweep & Sonnet & 1 & later (520) & 64 & 64 & 61 & 62 & 62 \\
fault-rate sweep & Sonnet & 4 & first (200) & 34 & 61 & 66 & 73 & 74 \\
fault-rate sweep & Sonnet & 4 & later (520) & 69 & 73 & 73 & 72 & 72 \\
\addlinespace
matched pairs & flash & 1 & first (100) & 18 & 30 & 32 & 29 & 29 \\
matched pairs & flash & 1 & later (352) & 38 & 43 & 43 & 43 & 43 \\
matched pairs & flash & 4 & first (100) & 35 & 64 & 71 & 72 & 74 \\
matched pairs & flash & 4 & later (352) & 36 & 38 & 38 & 37 & 37 \\
matched pairs, fault-free only & flash & 1 & first (50) & 32 & 56 & 56 & 56 & 56 \\
matched pairs, fault-free only & flash & 1 & later (174) & 36 & 39 & 40 & 39 & 39 \\
\addlinespace
matched pairs & terra & 1 & first (100) & 30 & 56 & 57 & 54 & 54 \\
matched pairs & terra & 1 & later (292) & 61 & 61 & 57 & 57 & 57 \\
matched pairs & terra & 4 & first (100) & 43 & 83 & 89 & 90 & 92 \\
matched pairs & terra & 4 & later (292) & 62 & 66 & 65 & 66 & 66 \\
matched pairs, fault-free only & terra & 1 & first (50) & 54 & 96 & 98 & 98 & 98 \\
matched pairs, fault-free only & terra & 1 & later (148) & 58 & 61 & 59 & 59 & 59 \\
\addlinespace
matched pairs & Sonnet & 1 & first (100) & 36 & 62 & 63 & 60 & 60 \\
matched pairs & Sonnet & 1 & later (246) & 60 & 50 & 47 & 46 & 47 \\
matched pairs & Sonnet & 4 & first (100) & 38 & 75 & 81 & 81 & 82 \\
matched pairs & Sonnet & 4 & later (246) & 63 & 61 & 59 & 61 & 61 \\
matched pairs, fault-free only & Sonnet & 1 & first (50) & 52 & 90 & 92 & 92 & 92 \\
matched pairs, fault-free only & Sonnet & 1 & later (127) & 59 & 54 & 53 & 52 & 52 \\
\bottomrule
\end{tabular}

\end{table}
